\chapter{Speculative IR, Deoptimization, and State Transfer}
\label{ch:speculative-state-transfer}

A speculative loop can use a narrow accumulator while every observed sum fits in a machine range.
When the next sum exceeds that range, returning a wrapped value would change the language.
Restarting the function can also be wrong if earlier iterations already produced output.
The runtime must resume a general computation at a precise continuation with the values and effects that continuation expects.
That requirement determines what information the optimizing IR must retain.

The earlier speculation chapter introduces guards and deoptimization.
Here we make the state-transfer contract explicit, execute it at every loop position in a finite test domain, and extend the same reasoning to virtual objects and representation migration.

\section{Represent assumptions and recovery together}
\label{sec:ch66-contract}

A \term{speculative IR} exposes assumptions, guards, and recovery states as part of optimization correctness.
An assumption is a proposition such as a value having the integer type or a class layout retaining a given version.
A guard establishes it along the fast path; invalidation may establish that previously compiled code can no longer rely on it.
Every failure path needs a destination that is semantically valid at that point.

Let $s_o$ be an optimized state and $s_b$ a baseline state at a designated continuation.
A reconstruction relation $R(s_o,s_b)$ says that their remaining executions produce the same required observations under the established assumptions.
The relation need not be a bijection: multiple optimized states can reconstruct to equivalent baseline states, and the optimized state may omit dead source variables.
This continuation-oriented view is central to formal treatments of dynamic deoptimization~\cite{fluckiger-2018-deoptimization}.

The recovery information must survive ordinary optimization.
If a value is used only by a deoptimization recipe, dead-code elimination must still regard that use as live or retain a sound way to recompute it.
Recomputation must itself be pure, terminating, and valid at the recovery point.
Re-executing an arbitrary source expression is not a safe substitute for a state map.

\section{Choose a continuation before effects commit}
\label{sec:ch66-loop}

The companion sums a sequence of exact integers and appends an event $(i,total)$ after each successful iteration.
A noninteger produces a type error at its index, preserving earlier events.
The baseline accumulator has arbitrary precision.
The fast accumulator accepts only next sums in $[-128,127]$; it guards before updating either the sum or the event sequence.

For input $[7,120,1]$, the first two sums are $7$ and $127$.
At index two the proposed sum is $128$, so the guard fails.
The reconstruction is the loop-header frame $(i,total)=(2,127)$ with the two already committed events.
The baseline executes only the remaining iteration and appends $(2,128)$.
Table~\ref{tab:ch66-trace} shows the complete trace.

\begin{table}[htbp]
\caption{Recovery resumes before the failing iteration and preserves the committed event prefix.}
\label{tab:ch66-trace}
\centering
\begin{tabular}{llll}
\toprule
Index & Incoming total & Decision & Newly committed event \\
\midrule
$0$ & $0$ & fast addition & $(0,7)$ \\
$1$ & $7$ & fast addition & $(1,127)$ \\
$2$ & $127$ & reconstruct, then baseline & $(2,128)$ \\
\bottomrule
\end{tabular}
\end{table}

Algorithm~\ref{alg:ch66-transfer} uses this single recovery boundary.
The type test must precede arithmetic on a value that may not support addition.
The range test is mathematical or uses a checked operation; performing an overflowing signed host addition before testing its range would not implement the contract.

\begin{algorithm}[htbp]
\caption{Run a narrow speculative loop with loop-header recovery.}
\label{alg:ch66-transfer}
\begin{pseudocode}
\Require input sequence, valid loop-header frame $(i,total,trace)$
\Ensure the baseline outcome and observable event sequence
\While{$i$ is below the input length}
  \State $v\gets$ input element at $i$
  \If{$v$ is not an exact integer}
    \State \Return baseline execution from the current frame
  \EndIf
  \State compute candidate sum without host-language overflow
  \If{the candidate is outside $[-128,127]$}
    \State reconstruct the current frame
    \State \Return baseline execution from that frame
  \EndIf
  \State commit the candidate sum and append its event
  \State $i\gets i+1$
\EndWhile
\State \Return successful baseline-format outcome
\end{pseudocode}
\end{algorithm}

Figure~\ref{fig:ch66-boundary} places the guard before the commit boundary.
If the guard were moved after output, recovery at the earlier loop header would duplicate output.
Moving it earlier can also fail if its reconstruction map references values not yet defined.
Guard motion is thus constrained by both assumptions and reconstructability.

\begin{figure}[htbp]
\centering
\begin{tikzpicture}[node distance=16mm and 18mm]
\node[basic block] (header) {header: i, total, trace};
\node[basic block,below=of header] (guard) {type and sum-range guard};
\node[basic block,below left=of guard] (fast) {commit sum and event};
\node[basic block,below right=of guard] (recover) {baseline at same header};
\draw[cfg edge] (header)--(guard);
\draw[cfg edge] (guard)--node[left,font=\small] {pass}(fast);
\draw[cfg edge] (guard)--node[right,font=\small] {fail}(recover);
\end{tikzpicture}
\caption{A failed guard transfers state without replaying effects from completed iterations.}
\label{fig:ch66-boundary}
\end{figure}

For $n$ elements the model performs $O(n)$ iterations, with arithmetic costs dependent on integer width.
Its diagnostic trace consumes $O(n)$ storage and is copied once on recovery; a real runtime usually stores committed external effects in the environment rather than copying them into a frame.
The trace is an oracle for this experiment, not a proposed production frame format.

\section{OSR is a second state-transfer problem}
\label{sec:ch66-osr}

\term{On-stack replacement} transfers an already active computation into another version, often at a loop header.
In the companion, a baseline prefix executes a chosen number of iterations, then the remaining loop enters the guarded implementation.
The transfer supplies index, accumulator, and the committed trace prefix.
The fast loop rechecks its own conditions, so a large incoming sum does not silently enter unchecked narrow arithmetic.

A native implementation that stores the incoming accumulator in an eight-bit location would first need an entry range guard, even if the next sum happened to fit.
Our Python model holds the incoming value exactly and checks every proposed update.
This illustrates why a representation-specific entry contract must be stated separately from the source-level continuation.

OSR and deoptimization need not be inverses at every machine instruction.
They can operate at selected safe boundaries, use different layouts, and transfer to different tiers.
Table~\ref{tab:ch66-recipes} records representative recovery sources.
Inlining may require reconstructing several logical frames from one optimized physical frame.

\begin{table}[htbp]
\caption{Recovery recipes explain how each required baseline value is obtained.}
\label{tab:ch66-recipes}
\centering
\begin{tabular}{ll}
\toprule
Recipe kind & Example obligation \\
\midrule
Location & read a live register or stack slot of known width \\
Constant & supply a value proved unchanged at this continuation \\
Pure recomputation & add two retained values without new effects \\
Virtual object & allocate and initialize the preserved object identity \\
Logical frame & restore the continuation of an inlined caller \\
\bottomrule
\end{tabular}
\end{table}

\section{Materialize identities before filling references}
\label{sec:ch66-objects}

Scalar replacement can eliminate an allocation from the fast path while retaining a recipe for the object's fields.
Recovery may then need to \term{materialize} the object.
Allocating a fresh object independently for each field reference breaks sharing.
Recursively building fields before recording the object breaks cycles.
Algorithm~\ref{alg:ch66-materialize} solves both problems with two passes.

\begin{algorithm}[htbp]
\caption{Materialize a finite virtual-object graph while preserving identity.}
\label{alg:ch66-materialize}
\begin{pseudocode}
\Require well-formed recipes with stable object IDs and valid field references
\Ensure concrete objects with the recipe graph's sharing and cycles
\ForAll{virtual object IDs $v$}
  \State allocate an uninitialized shell and record it as the image of $v$
\EndFor
\ForAll{virtual object IDs $v$}
  \ForAll{fields in the recipe for $v$}
    \State write the constant value or already allocated referenced shell
  \EndFor
\EndFor
\State expose the completed objects to the recovered frame
\end{pseudocode}
\end{algorithm}

For $V$ virtual objects and $F$ fields, the graph work is $O(V+F)$ with expected constant-time identity lookup, excluding allocation and collector overhead.
Figure~\ref{fig:ch66-cycle} shows a two-object cycle.
The companion checks that following two \code{next} fields returns the identical first object, not merely an equal-looking copy.
Production implementations must also root intermediate shells for the collector, handle allocation failure according to the language contract, and prevent observation of incomplete objects.

\begin{figure}[htbp]
\centering
\begin{tikzpicture}[node distance=28mm]
\node[basic block] (a) {object a};
\node[basic block,right=of a] (b) {object b};
\draw[cfg edge,bend left=25] (a) to node[above,font=\small] {next}(b);
\draw[cfg edge,bend left=25] (b) to node[below,font=\small] {next}(a);
\end{tikzpicture}
\caption{Allocating both shells first lets recovery preserve a cyclic object graph.}
\label{fig:ch66-cycle}
\end{figure}

\section{Verification and exercises}
\label{sec:ch66-exercises}

The tests enumerate every sequence of length zero through four over six values: $-129$, $-1$, $0$, $7$, $127$, and a noninteger.
Each sequence is tested at every possible OSR prefix position.
They compare final status, error index, exact accumulator, and the whole event trace against baseline execution.
This detects duplicated effects, skipped inputs, wrong error locations, and lost values in the finite domain.
It does not prove correctness of a native stack-map format or garbage collector.

\begin{enumerate}
\item\label{ex:ch66-overflow} Work the guarded execution of $[7,120,1]$ and give the recovered frame and final event trace.
\item\label{ex:ch66-error} Execute $[7,\text{noninteger},4]$ with OSR before index one.
Give the error position, accumulator, and committed events.
\item\label{ex:ch66-replay} Explain the observable bug if recovery restarts the function after the first two events of the overflow example.
\item\label{ex:ch66-identity} Two virtual objects point to each other and a recovered local aliases the first.
Describe a two-pass recipe and the identity assertions that distinguish it from deep copying.
\end{enumerate}

\section{Further reading}
\label{sec:ch66-reading}

Flückiger and colleagues formalize correctness obligations for speculative optimization and dynamic deoptimization~\cite{fluckiger-2018-deoptimization}.
Chapter~\ref{ch:continuous-compilation} uses these transfer contracts to manage multiple live code and data representations.
